\section{Introduction}
\label{sec:introduction}

The capacity for continuous, sequential adaptation without the eradication of
historical knowledge remains an architectural cornerstone of biological
intelligence, yet it presents a fundamental barrier for artificial neural
networks. In deep learning paradigms, sequential training across non-stationary
data distributions triggers a systemic pathology known as \emph{catastrophic
forgetting} \citep{mccloskey_catastrophic_1989}. When a network optimizes its
parameters for an incoming task, gradient updates blindly overwrite the weight
configurations essential for retaining prior task spaces, shifting the internal
representation manifolds away from historical objective minima.

Modern continual learning (CL) literature is conventionally structured into
three core algorithmic families: replay-based strategies that leverage episodic
data buffers, regularization-based frameworks that penalize parameter deviations,
and architectural or parameter-isolation methods \citep{de_lange_continual_2022}.
Within the meta-learning domain, hypernetworks \citep{ha_hypernetworks_2016}
have emerged as an elegant paradigm. Hypernetworks circumvent parameter-space
conflicts by utilizing a centralized, lightweight meta-model to generate the
target network's complete parameter weights conditioned on task embeddings, as
formalized by \citet{oswald_continual_2022}. Concurrently, within the field of
geometric optimization, gradient projection methods---exemplified by the
recently proposed Fisher-Orthogonal Projected Natural Gradient Descent (FOPNG)
\citep{garg_fisher-orthogonal_2026}---explicitly constrain subsequent updates to
the null-space of historical parameter spaces.

Despite the mathematical elegance of both hypernetworks and gradient projection
methods, their structural intersection uncovers a critical and previously
unexamined optimization pathology. To prevent a catastrophic parameter explosion
within the meta-model, a hypernetwork must use sequential weight-chunking.
Depending on the dimensional footprint of the target architecture, this
constraint forces the target model to be produced in hundreds or thousands of
individual sub-blocks. Executing this spawning mechanism prior to a single
backward pass accumulates gradients additively, resulting in parameter gradients
of exceptionally high magnitude. This explosion numerically erases the stabilizing
effect of $\lambda$ during matrix inversion, inflating the condition number of
the projection kernel to extreme scales (e.g., exceeding $10^{10}$) and inducing
catastrophic projection failures.

\subsection*{Research Objectives and Algorithmic Framework}

The primary research objective of this thesis is to systematically resolve this
linear-algebraic collapse, map the initialization mechanics of meta-projection
pipelines, and evaluate whether the integration of hard geometric constraints on
top of soft functional regularizers yields an optimized sequential learning system.

To achieve this, we introduce projection-scoped normalization (column-wise
$\ell_2$ orthonormalization of $G$ combined with absolute-maximum scaling of
$\hat{F}$), an \texttt{AdamW} first-task initialization protocol, and
\texttt{iFOPNG}---an inertial FOPNG variant using the combined metric
$F_c = \hat{F}_{\mathrm{new}} + \hat{F}_{\mathrm{old}}$.  These contributions
enable stable gradient projection in chunked hypernetworks, with direct
applicability to on-device sequential learning where retaining adaptive momentum
is critical.

\subsection*{Research Questions}

To systematically navigate these boundaries, this thesis addresses the following
primary research question:
\begin{quote}
\emph{To what extent can gradient projection methods be successfully integrated
into chunked hypernetwork architectures for continual learning, and how do their
structural pathologies, initialization protocols, and inertial variants impact
optimization efficiency?}
\end{quote}
This is decomposed into four sub-research questions (Sub-RQs):
\begin{itemize}
    \item[\textbf{Sub-RQ1}] \textbf{(Synergy vs.\ Conflict):} Does the joint integration
    of gradient projection frameworks and soft output-scoped functional
    regularization yield a constructive synergy, or does a standalone regularized
    hypernetwork offer a more efficient optimization frontier?
    \item[\textbf{Sub-RQ2}] \textbf{(Gradient Pathology):} How does the architectural
    necessity of spawning thousands of sequential weight-chunks impact gradient
    magnitudes, and to what extent can local projection-scoped normalization
    prevent the resulting explosion from collapsing the matrix inversion?
    \item[\textbf{Sub-RQ3}] \textbf{(Parameter Inertia):} To what degree does
    evaluating the natural gradient projection under a joint Riemannian metric
    ($F_c = \hat{F}_{\mathrm{new}} + \hat{F}_{\mathrm{old}}$), rather than the task-specific
    metric ($\hat{F}_{\mathrm{new}}$) used in standard FOPNG, influence task retention?
    \item[\textbf{Sub-RQ4}] \textbf{(Fisher Accumulation):} How should the historical
    component $\hat{F}_{\mathrm{old}}$ be accumulated across tasks? Specifically, does
    a non-decaying maximum operator preserve early-task knowledge more
    effectively than an exponential moving average (EMA)?
\end{itemize}

\subsection*{Summary of Empirical Findings}

Projection-scoped normalization completely resolves the matrix inversion
collapse; \texttt{iFOPNG} consistently outperforms \texttt{FOPNG} wherever
inter-task Fisher overlap is substantial; and the MAX accumulation operator is
the preferred hyperparameter-free default.  Critically, normalized projection
methods nonetheless fail to match a hypernetwork optimized via \texttt{Adam}:
the loss of adaptive momentum proves to be the binding constraint in compressed
generative weight spaces.